%%%%%%%%%%%%%%%% Pakete & Klassen %%%%%%%%%%%%%%%%%%%%%%%%%%%%%%%%%%%%%%%%%%%%%%%%%%%%%%%%%%%%%%%%%%%%%%%%%%%%%%%%%%%%%%%%%%%%%%%%%%%%
\documentclass[a4paper, 12pt, oneside]{article}             % Papierformat, Schriftgröße, einseitiges Layout; Verwendung des KOMA-Script
%\usepackage[ansinew, latin1]{inputenc}                     % Zeichenkodierung: Übersetzung von Umlauten, Windows/Unix-Umgebung; deutsche Umlaute
\usepackage[utf8]{inputenc}                                 % Zeichenkodierung
\usepackage[T1]{fontenc}                                    % Schriftart
\usepackage[ngerman]{babel}                                 % neue deutsche Rechtscheibung, Silbentrennung
\usepackage{csquotes}                                       % "Schlaue" Anführungszeichen
\usepackage{lmodern}                                        % Schriftart
\usepackage{textcomp}                                       % führt Sonderzeichen ein
\usepackage{url}                                            % URLs sind Hyperlinks
\usepackage{hyperref}                                       % Verlinkung, jeder Querverweis kann im PDF angeklickt werden
\usepackage{color}                                          % Farben ermöglichen. Grundfarben: black, white, red, green, yellow, blue
\usepackage{longtable}                                      % für extra lange Tabellen
\usepackage{enumitem}                                       % schönere Aufzählungen
\usepackage{multirow}                                       % \multirow{} und \multicolumn{} für Tabellen
\usepackage{graphicx}                                       % erweitertes Paket für Bilder-Einbindung (EPS, PNG, JPG, PDF)
\usepackage[left=1.8cm,right=1.8cm,top=1.8cm]{geometry}     % Seitenränder
\usepackage[final]{pdfpages}				                % Einfügen von mehrseitigen PDF-Dateien,
						  % benutzen mit \includepdf[pages=-, pagecommand={\thispagestyle{headings}}]{doc.pdf}
\pagestyle{myheadings}                                      % ermöglicht Kopfzeilen
\MakeOuterQuote{"}                                          % Anführungszeichen automatisch umwandeln


%%%%%%%% Sachen definieren %%%%%%%%%%%%%%%%%%%%%%%%%%%%%%%%%%%%%%%%%%%%%%%%%%%%%%%%%%%%%%%%%%%%%%%%%%%%%%%%%%%%%

\makeatletter
\newcounter{protokollcounter}
\newcommand{\protokollparse}[1]{\@for\@tempa:=#1\do{\stepcounter{protokollcounter}}}

\newcommand{\protokolldatum}{00:00:0000}
\newcommand{\protokollbeginn}{00:00 Uhr}
\newcommand{\protokollende}{00:00 Uhr}
\newcommand{\protokollant}{NN}
\newcommand{\protokollanwesend}{NN}
\newcommand{\protokollfehlend}{NN}
\newcommand{\cntprotokollanwesend}{\setcounter{protokollcounter}{0}\protokollparse{\protokollanwesend}\arabic{protokollcounter}}
\newcommand{\cntprotokollfehlend}{\setcounter{protokollcounter}{0}\protokollparse{\protokollfehlend}\arabic{protokollcounter}}

\newcommand{\protokolltitlepage}{%
  \begin{center}
        \includegraphics[width=0.75\textwidth]{fs-physik-logo.pdf}\\
        \vspace{0.3cm}
        \Huge{\textbf{Protokoll der Fachschaftssitzung} \\
        \vspace{0.2cm}
        \protokolldatum\\
        \vspace{0.2cm}
        }
        \normalsize

		\begin{longtable}{|p{6.1cm}|p{6.1cm}|}
			\hline
			\textbf{Beginn:}    \protokollbeginn &
			\textbf{Ende:}      \protokollende \\
			\hline
			\textbf{Protokollant:} &
			\protokollant \\
			\hline
			\multicolumn{2}{|p{12.5cm}|}{\textbf{Anwesende (\cntprotokollanwesend): }} \\
			\multicolumn{2}{|p{12.5cm}|}{\protokollanwesend}\\
			\hline
			\multicolumn{2}{|p{12.5cm}|}{\textbf{Entschuldigt (\cntprotokollfehlend): }} \\
			\multicolumn{2}{|p{12.5cm}|}{\protokollfehlend}\\
			\hline
		\end{longtable}
	\end{center}
	\tableofcontents
}


%%% Datum, Uhrzeit, Teilnehmer, Ort
\renewcommand{\protokolldatum}{dd.mm.yyyy}
\renewcommand{\protokollbeginn}{hh:mm Uhr}
\renewcommand{\protokollende}{hh:mm Uhr}
\renewcommand{\protokollant}{Protokollant}
\renewcommand{\protokollort}{StudiO, Wilhelm-Klemm-Str.\ 9, 48149 Münster}
% Listen unbedingt durch Kommata trennen, sonst wird nicht richtig gezählt.
\renewcommand{\protokollanwesend}{Namen}
\renewcommand{\protokollweitere}{Namen}

%%% Beginn des Dokuments
\begin{document}
% Titelseite
\protokolltitlepagekonstituierend

%%% Dokument-Inhalt
\section{Begrüßung und Wahl einer Versammlungsleitung und eines Protokollanten}
\begin{itemize}
	\item XXX begrüßt alle Anwesenden.
	\item YYY wird mit 00:00:00 (J:N:E)\footnote{Hinweis zur Konvention: Ergebnisse von Abstimmungen werden im in der Reihenfolge (Ja-Stimmen:Nein-Stimmen:Enthaltungen) angegeben.} zur Versammlungsleitung gewählt.
	\item ZZZ wird mit 00:00:00 zum Protokollanten gewählt.
\end{itemize}

\section{Feststellung der Beschlussfähigkeit}
Für die Beschlussfähigkeit müssen mindestens \SI{50}{\percent} der wahlberechtigten Mitglieder anwesend sein.

Die Beschlussfähigkeit wird mit XX anwesenden von insgesamt 15 Mitgliedern der FSV festgestellt.

\section{Feststellung der Tagesordnung}
Festgestellt wie vorgestellt.

\section{Wahl des neuen Fachschaftsrats}
\begin{tabular}{| r @{ } l | p{5cm} | p{3cm} |}
\hline
& "Amt" (kann entfallen)
& \centering\let\newline\\\arraybackslash Vor- und Nachname;\newline Matrikel-Nr.
& \centering\let\newline\\\arraybackslash Abstimmung\newline(Ja:Nein:Enth.)
\\ \hline\hline
1.  & Vorsitzende/r & &
\\ \hline
2.  & Stellv. Vorsitzende/r & &
\\ \hline
3.  & Finanzreferent/in & &
\\ \hline
4.  & Stellv. Finanzreferent/in & &
\\ \hline
5.  & Erstibeauftragte/r & &
\\ \hline
6.  & Beauftragte/r O-Woche & &
\\ \hline
7.  & Stellv. Beauftragte/r O-Woche & &
\\ \hline
8.  & Beauftragte/r Ersti-Fahrt & &
\\ \hline
9.  & Stellv. Beauftragte/r Ersti-Fahrt & &
\\ \hline
10. & Beauftragte/r Ersti-Fibel & &
\\ \hline
11. & Stellv. Beauftragte/r Ersti-Fibel & &
\\ \hline
12. & Evaluationsbeauftragte/r & &
\\ \hline
13. & Stellv. Evaluationsbeauftragte/r & &
\\ \hline
14. & Stellv. Evaluationsbeauftragte/r & &
\\ \hline
15. & Beauftragte/r Sommerfest & &
\\ \hline
16. & Stellv. Beauftragte/r Sommerfest & &
\\ \hline
17. & Stellv. Beauftragte/r Sommerfest & &
\\ \hline
18. & Eventmanager/in & &
\\ \hline
19. & Stellv. Eventmanager/in & &
\\ \hline
20. & E-Mail-Beauftragte/r & &
\\ \hline
21. & Beauftragte/r BaMa-Tag & &
\\ \hline
22. & Stellv. Beauftragte/r BaMa-Tag & &
\\ \hline
23. & Beauftragte/r Buchmarkt & &
\\ \hline
24. & Stellv. Beauftragte/r Buchmarkt & &
\\ \hline
25. & Beauftragte/r ZaPF/Vernetzung & &
\\ \hline
26. & Beauftragte/r E-Learning & &
\\ \hline
27. & Beauftragte/r Protokoll-/Klausurverleih & &
\\ \hline
28. & Beauftragte/r Protokoll-/Klausurverleih & &
\\ \hline
29. & Pizzabeauftragte/r & &
\\ \hline
\end{tabular}

\bigskip
Das Protokoll wurde genehmigt auf der Sitzung am \protokolldatum.

\bigskip
\bigskip
Unterschrift der Protokollführerin/des Protokollführers

\end{document}
